% Generated draft: requires mathematical and source review.
\documentclass[11pt]{article}
\usepackage[margin=25mm]{geometry}
\usepackage{fontspec}
\setmainfont{Latin Modern Roman}
\usepackage{amsmath,amssymb,mathtools,unicode-math}
\setmathfont{Latin Modern Math}
\usepackage{enumitem,needspace,xurl,hyperref}
\setlength{\parindent}{0pt}
\setlength{\parskip}{5pt}
\setlength{\emergencystretch}{3em}
\newcommand{\N}{\mathbb{N}}
\newcommand{\Z}{\mathbb{Z}}
\newcommand{\Q}{\mathbb{Q}}
\newcommand{\R}{\mathbb{R}}
\newcommand{\C}{\mathbb{C}}
\newcommand{\F}{\mathbb{F}}
\newcommand{\A}{\mathbb{A}}
\newcommand{\PP}{\mathbb{P}}
\newcommand{\E}{\mathbb{E}}
\newcommand{\eps}{\varepsilon}
\newcommand{\dd}{\,\mathrm{d}}
\begin{document}
\section*{UnsolvedMath ID: 20001151}
\textbf{Source problem code:} catalogue entry
\subsection{Definitions and mathematical statement}
Let \(G\) be a nonempty finite simple graph. An even hole is an induced cycle of even length at least four. Let \(\chi(G)\) denote chromatic number and \(\omega(G)\) clique number. The source's conjecture, attributed to Ho\`ang, is
\[
 G\text{ even-hole-free}\quad\Longrightarrow\quad
 \chi(G)\le2\omega(G)-1.
\]
Nonemptiness is explicit here because the literal right side is negative for the null graph under \(\omega(\varnothing)=\chi(\varnothing)=0\). The appended sentence ``Hayward and Reed proposed the following'' introduces subsequent source material and does not supply another assertion in this record.

The supplied aid reports the conjecture resolved by a 2023 theorem of Chudnovsky and Seymour: every non-null even-hole-free graph has a bisimplicial vertex, meaning its neighborhood can be covered by two cliques. Since induced subgraphs retain even-hole-freeness, successive choices give an ordering \(v_1,\ldots,v_n\) with \(G_i=G[\{v_i,\ldots,v_n\}]\) and
\[
 |N_{G_i}(v_i)|\le2\omega(G_i)-2.
\]
The aid states the stronger adaptive list-coloring consequence: for arbitrary finite color lists satisfying
\[
 |L(v_i)|\ge2\omega(G_i)-1\quad(1\le i\le n),
\]
a proper list coloring exists by reverse greedy coloring in this ordering. In particular
\[
 \chi(G)\le\chi_\ell(G)\le2\omega(G)-1,
\]
where \(\chi_\ell\) is the least uniform list size guaranteeing a proper coloring for every assignment of lists of that size. Odd holes \(C_{2t+1}\), \(t\ge2\), give sharpness at \(\omega=2\); this is not a claim of sharp examples for every clique number.

The aid's separate conditional equality lemma is: if an even-hole-free graph \(G\) has \(\omega(G)=w\ge2\) and \(\chi(G)=2w-1\), and \(H\) is any induced-minimal subgraph with \(\chi(H)=\chi(G)\), then \(\omega(H)=w\), and there is a vertex \(v\in V(H)\) such that
\[
 d_H(v)=2w-2,\qquad
 N_H(v)=Q_1\mathbin{\dot\cup}Q_2,
 \qquad |Q_1|=|Q_2|=w-1,
\]
where each \(Q_i\) is a clique, and
\[
 \chi(H-v)=2w-2,\qquad\omega(H-v)=w.
\]
The two cliques are not asserted anticomplete. ``Induced-minimal'' means no proper induced subgraph of \(H\) has that same chromatic number. This lemma is conditional on an extremal graph's existence and is not a classification or an existence assertion for all \(w\). All theorem, list, and equality conclusions are attributed to the supplied aid and require citation/proof verification.
\subsection{Short English statement}
The source asks for the bound \(\chi\le2\omega-1\) for even-hole-free graphs. The supplied aid reports the bisimplicial-vertex theorem as its resolution, strengthens the bound to list coloring, and gives a conditional local structure lemma for chromatically extremal graphs.
\subsection{Sources}
\begin{itemize}
\item Local UnsolvedMath registry, permanent ID \(20001151\), source code \texttt{catalogue entry}; \texttt{aim-combinatorics-notes.json}, \texttt{notes[275]}, tag \texttt{conjecture}.
\item AIM, \emph{The Perfect Graph Conjecture}, official PDF version dated 24 August 2004, Chapter G.4, ``Even-Hole-Free Graphs,'' Conjecture 3, Ho\`ang: \url{https://aimath.org/WWN/perfectgraph/perfectgraph.pdf}. Statement status \texttt{exact}. The trailing Hayward--Reed transition is retained as context only. No PDF was accessed for this draft.
\item Chudnovsky and Seymour, bisimplicial-vertex theorem reported as proved in 2023. No full title, theorem number, or URL is supplied. The adaptive list consequence and candidate equality lemma have no separate authoritative citation in the record.
\item Research attempt 1: \texttt{full\_solution}, status \texttt{solved}, classification \texttt{SOLVED-IN-LITERATURE}, rigor label \texttt{proved}, novelty confidence low. Created and updated 14 August 2026. No original remarks or literature field are supplied. This machine-generated aid is not peer reviewed; independently verify before citation.
\item Ulam AI / UnsolvedMath Contributors; CC BY 4.0.
\end{itemize}
\end{document}
